\section{Limites et ouvertures}
\label{sec:limites_ouvertures}

\subsection{Ce qui est solidement acquis}
\label{sec:acquis}

Le programme HP-PT, dans son état actuel, fournit~:
\begin{itemize}[nosep,leftmargin=*]
\item l'expression close des amplitudes $\kappap(s)$ et la
représentation Fredholm exacte de $\Rres(s)$
(thm.~\ref{thm:fredholm}) ;
\item l'opérateur canalisé $\Hilb_{\mathrm{PT}}$ pour le couplage avec
$\epsstar^{\mathrm{geom}}$ canonique sans paramètre libre
(lem.~\ref{lem:epsstar_geom}) ;
\item le cut-off dynamique $\umax(\gamma) = \gamma/\sqrt{2\pi}$ forcé
par la cellule de Planck symplectique
(prop.~\ref{prop:cellule_planck}) ;
\item la densité semi-classique exacte
$\NPT(\gamma) = N_{\mathrm{RvM}}(\gamma) + 1/8$ (forme fonctionnelle
identique) (prop.~\ref{prop:densite_NPT}) ;
\item les indices de défaut $(1, 1)$ et la condition aux bords
antipériodique forcée par $T_3$ (prop.~\ref{prop:antiperiodicite}) ;
\item la formule de trace analytique exacte $-d \log \Rres/ds$
(thm.~\ref{thm:trace_formula}) ;
\item la trace régularisée opératorielle $\Treg(s)$ avec pôles
distributionnels aux zéros de $\Rres$ (\S\ref{sec:trace_reg}) ;
\item la triple identification cohomologique du résidu $1/8$
(\eqref{eq:triple_identification}) ;
\item la dualité structurelle Higgs $\leftrightarrow \zeta$
(conjecture~\ref{conj:K4}).
\end{itemize}

\paragraph{Tableau récapitulatif des statuts épistémiques.}

\begin{longtable}{@{}p{8cm}p{2.5cm}p{4cm}@{}}
\toprule
\textbf{Résultat} & \textbf{Statut} & \textbf{Source}\\
\midrule
\endhead
$\sper = 1/2$ ($\Tone$) & \textbf{[Thm]} & \cite{PT_M1}\\
$\det(I - \DRop)^{-1} = \Rres(s)$, $\Reop(s) > 1$ & \textbf{[Thm]} & thm.~\ref{thm:fredholm}\\
Formule de trace $-d\log \Rres/ds$ dans $\Reop(s) > 1$ & \textbf{[Thm]} & thm.~\ref{thm:trace_formula}\\
Indices de défaut $(1,1)$ et BC AP $\theta = \pi$ & \textbf{[Thm]} & prop.~\ref{prop:antiperiodicite}\\
$c_1 = 1/2$ sur fibré spinoriel $\qplus/\qminus$ & \textbf{[Thm]} & thm.~\ref{thm:berry_3pi}\\
Densité $\NPT = N_{\mathrm{RvM}} + 1/8$ (identité asymptotique) & \textbf{[Id]} & prop.~\ref{prop:densite_NPT}\\
$1/8 = c_1/N_{\mathrm{corners}}$ (rigueur APS stricte) & \textbf{[Thm partiel]} & prop.~\ref{prop:c1_sur_N}\\
$\zetaplus\,\zetaminus \ne 0$ sur ligne critique & \textbf{[CONJ]} & rem.~\ref{rem:non_annulation}\\
Prolongement analytique strict de $\Rres(s)$ hors $\Reop(s) > 1$ & \textbf{[OPEN]} & ouv.~1 ci-dessous\\
Pôles distributionnels de $\Treg(s)$ sur $\Reop(s) = 1/2$ uniquement & \textbf{[OPEN]} ($\equiv$~RH) & ouv.~2 ci-dessous\\
Choix canonique $u_\star$ ($\sqrt{2\pi}$ vs $\log 3$) & \textbf{[OPEN]} & ouv.~4 ci-dessous\\
Dualité Higgs $\leftrightarrow \zeta$ structurelle (K4) & \textbf{[CONJ forte]} & conj.~\ref{conj:K4}\\
$\Delta_N(L \bmod q) = 1/(8q)$ pour Dirichlet (P3) & \textbf{[FALSIFIED, $T \le 200$]} & \S\ref{sec:M1bis}\\
$G_{\mathrm{HP4.d}}$~: spectre direct = $\{\gamma_n\}$ & \textbf{[FALSIFIED]} & \S\ref{sec:trace_reg}\\
\bottomrule
\caption{Statuts épistémiques des résultats du programme.}
\label{tab:statuts}
\end{longtable}

\subsection{Ce qui reste ouvert}
\label{sec:ouvertures}

Quatre ouvertures structurelles, plus une hypothèse de non-annulation
sous-jacente.

\paragraph{Ouverture 1 (G\_HP3.a).}
Prolongement analytique de $\Rres(s)$ hors $\Reop(s) > 1$. La trace
$\Treg(s)$ est définie par série dans $\Reop(s) > 1$ puis prolongée
distributionnellement à $\Reop(s) = 1/2$ ; la région intermédiaire
$1/2 < \Reop(s) < 1$ demande l'analyse classique de $\Rres$. C'est
essentiellement le problème classique du prolongement de $\zeta$.

\paragraph{Ouverture 2 (RH stricte).}
Prouver que \textbf{tous} les pôles distributionnels de $\Treg(s)$ sont
sur $\Reop(s) = 1/2$. Le programme établit que la structure favorise
cette ligne (quatre mécanismes équivalents,
\S\ref{sec:localisation_critique}), mais ne prouve pas l'exclusion
stricte des pôles hors ligne. C'est essentiellement RH.

\paragraph{Ouverture 3 (K3 strict APS).}
Rigueur formelle Atiyah--Patodi--Singer pour la variété à coin de la
double barrière BK, au-delà du calcul d'aire et de la vérification
numérique. Le mécanisme cohomologique est identifié
(\S\ref{sec:indice_APS}), sa formalisation stricte reste à développer.

\paragraph{Ouverture 4 (K1, choix canonique de $u_\star$).}
La barrière inférieure $u_\star$ admet deux candidats naturels~:
$u_\star = \sqrt{2\pi}$ (point fixe de l'involution
$u \leftrightarrow p_u$, choix retenu ici) et $u_\star = \log 3$
(premier prime actif PT, choix arithmétique). Les deux produisent la
même densité asymptotique $N(\gamma)$ mais des corrections
$O(1/\gamma)$ différentes. Le statut canonique du choix entre les deux
reste ouvert. L'article retient $u_\star = \sqrt{2\pi}$ par symétrie
d'involution, sans prétendre clore K1.

\paragraph{Hypothèse \textbf{[CONJ]} sous-jacente : non-annulation $\zetaplus\,\zetaminus \ne 0$.}
L'identification entre zéros de $\Rres$ et zéros de $\zeta$ utilisée
tout au long de l'article repose sur cette non-annulation. Elle est
vérifiée numériquement sur les zones testées mais n'est pas démontrée
(rem.~\ref{rem:non_annulation}). Un échec partiel introduirait des zéros
supplémentaires (zéros de $\zeta_\pm$) qui se traduiraient en pôles
supplémentaires de $\Treg(s)$ ; ces zéros possibles doivent être
traités comme conjecturalement absents.

\subsection{Position canonique du programme}
\label{sec:position_canonique}

Le programme HP-PT est un cadre opératoriel précis pour les zéros de
$\zeta$. Il fournit~:
\begin{enumerate}[nosep,leftmargin=*]
\item une expression close des zéros via $\kappap(s)$ (intermédiaire
$\DRop$) ;
\item une formule de trace régularisée canonique $\Treg(s)$ ;
\item quatre mécanismes équivalents pour $\Reop(s) = 1/2$ ;
\item un dictionnaire continu/discret unifié par la persistance ;
\item une explication structurelle des coïncidences numériques connues
($2\pi$, $7/8$, $1/8$, ligne critique).
\end{enumerate}

Il n'apporte pas~:
\begin{enumerate}[nosep,leftmargin=*]
\item une preuve stricte de RH classique ;
\item un raccourci à l'analyse distributionnelle de la trace
régularisée.
\end{enumerate}

Cette distinction est maintenue tout au long du programme~: on cherche à
\textbf{comprendre} la structure des zéros, pas à les \textbf{prouver}
sur la ligne critique. Le passage à la preuve formelle demande
l'analyse de Schwartz fine sur $\Treg(s)$, qui n'est pas plus facile que
la preuve classique de RH.

\subsection{Comparaison avec Berry--Keating--Connes}
\label{sec:comparaison_BKC}

\begin{table}[ht]
\centering
\small
\begin{tabular}{@{}p{3cm}p{4.5cm}p{5cm}@{}}
\toprule
\textbf{Aspect} & \textbf{Berry--Keating--Connes} & \textbf{PT}\\
\midrule
Espace de phases & $(x, p) \in \Reals_+^{2}$ & $(u, p_u)$, $u = \log p$\\
Opérateur & $H = xp + 1/2$ & $\HPTBK = (u\,p_u + p_u\,u)/2$\\
Mesure & $dx\,dp / 2\pi$ & $du\,dp_u / 2\pi$\\
Cut-off & $x \ge \ell_x$, $p \ge \ell_p$, $\ell_x \ell_p = 2\pi$ & $u_\star\,p_\star = 2\pi$ (cellule de Planck PT)\\
Régularisation & adélique (Connes) & atomique PT\\
Trace & $\Tr(K) = \sum_\rho h(\gamma_\rho)$ (Weil) & $\Tr_{\mathrm{reg}} = \sum_p \kappap [\text{Euler + couplage}]$\\
Couplage holonomique & aucun (zêta seule) & $\kappap$ (bifurcation $\qplus/\qminus$)\\
$2\pi$ & non interprété structurellement & $[u, p_u] = i$, Liouville canonique\\
$1/8$ & phase Maslov des coins (artefact) & $\sper^{2}/2 = c_1/4 = \sper/4$ (cohomologie spinorielle)\\
Spin $\sper = 1/2$ & postulé & sur-déterminé par $\sper^{2}/2 = \sper/4$ modulo $\Tone$\\
\bottomrule
\end{tabular}
\caption{Comparaison Berry--Keating--Connes vs PT~: l'apport principal de
PT est l'enrichissement intrinsèque (aucun nombre n'est posé ad hoc).}
\label{tab:comparaison_BKC}
\end{table}

\paragraph{Précisions historiques.}
L'opérateur originel de \citeauthor{BerryKeating1999} (1999) est
$H = xp$ ; la forme $H = xp + 1/2$ apparaît dans la formalisation
ultérieure de \citeauthor{Sierra2008} (2008) et Sierra--Townsend, qui
régularise $xp$ via une cellule de Planck et exhibe l'ordre symétrique.
La régularisation de \citeauthor{Connes1996} (\citeyear{Connes1996},
\citeyear{Connes1999}) est de nature différente~: elle s'effectue sur
l'espace adélique des classes d'idèles, par formule de trace
non-commutative. La formulation PT est atomique sur les premiers~: ni
adélique, ni continue sur $\Reals_+$.

\subsection{Perspectives}
\label{sec:perspectives}

Trois prolongements naturels.

\begin{itemize}[nosep,leftmargin=*]
\item \textbf{Autres fonctions $L$.} La prédiction Dirichlet ayant été
falsifiée, l'extension aux fonctions $L$ doit être repensée. La
cohomologie du fibré spinoriel sous \emph{twist} par $\chi$ modifie
complètement $c_1$ ; une reformulation cohomologique propre est
ouverte.
\item \textbf{Phase Maslov dans la phase de Berry continue.}
L'hypothèse propose que le $1/(8q)$ existe au niveau de la phase de
Berry continue le long du flot spectral, mais ne se manifeste pas dans
le compte des zéros. À tester.
\item \textbf{Rigueur APS stricte.} Développer le théorème d'indice
formel sur la variété à coin de la double barrière (K3 strict). Pas
urgent pour le cadre opératoriel, important pour la cohérence
structurelle.
\end{itemize}

